\documentclass[10pt,twocolumn]{article}

% ============================================================
% PACOTES
% ============================================================
\usepackage[a4paper,top=2cm,bottom=2.2cm,left=1.6cm,right=1.6cm,columnsep=0.6cm]{geometry}
\usepackage{times}
\usepackage[T1]{fontenc}
\usepackage[utf8]{inputenc}
\usepackage{amsmath,amssymb}
\usepackage{graphicx}
\usepackage{caption}
\usepackage{titlesec}
\usepackage{fancyhdr}
\usepackage{lastpage}
\usepackage{enumitem}
\usepackage{hyperref}
\usepackage{indentfirst}  % <--- ADICIONADO: Recua primeiro parágrafo

\hypersetup{colorlinks=true,linkcolor=blue,citecolor=blue,urlcolor=blue}

% ---------- Cabeçalho e rodapé ----------
\pagestyle{fancy}
\fancyhf{}
\fancyhead[L]{\footnotesize Universiade Estadual do Maranhão (UEMA)}
\renewcommand{\headrulewidth}{0.4pt}
\renewcommand{\footrulewidth}{0pt}

% ---------- Formatação de seções ----------
\titleformat{\section}
{\normalfont\centering\bfseries}
{\Roman{section}.}{0.6em}{\MakeUppercase}
\titleformat{\subsection}
{\normalfont\itshape}
{\Alph{subsection}.}{0.6em}{}

% ---------- Recuo de parágrafos ----------
\setlength{\parindent}{1.5em}  % <--- AUMENTEI PARA MELHOR VISIBILIDADE
\setlength{\parskip}{0pt}

% ---------- BIBER / BIBLATEX ----------
\usepackage[
backend=biber,           % Usa Biber como backend
style=numeric,           % Estilo numérico [1], [2], [3]...
% style=ieee,            % Alternativa: estilo IEEE
% style=abnt,           % Alternativa: estilo ABNT (requer pacote abntex2)
sorting=none,            % Ordem de citação (none = aparece na ordem citada)
% sorting=ynt,          % Ordena por ano, nome, título
maxbibnames=99,          % Mostra todos os autores na bibliografia
minbibnames=1,           % Mínimo de autores antes de "et al."
url=true,                % Inclui URLs
doi=true,                % Inclui DOIs
isbn=true,               % Inclui ISBNs
backref=true,            % Mostra páginas onde a referência foi citada
backrefstyle=three,      % Estilo dos backreferences
]{biblatex}

% Adiciona arquivo .bib
\addbibresource{referencias.bib}  % Nome do seu arquivo de referências

\hypersetup{colorlinks=true,linkcolor=blue,citecolor=blue,urlcolor=blue}

\begin{document}
	
	% ============================================================
	% TÍTULO E AUTORES
	% ============================================================
	\twocolumn[{%
		\centering
		{\LARGE\bfseries Pré-trabalho de Vibrações Mecânicas: Peteca no uso de análise da integridade estrutural de pisos}\par\vspace{1em}
		{\normalsize {Aurélio B. L. do Nascimento; Helder F. S. Nogueira; João G. M. Silva.}}\par\vspace{1.2em}
		\begin{minipage}{0.94\linewidth}
			\small
			%\textbf{Abstract}--- ...
			%\medskip
			%\textbf{Index Terms}--- ...
		\end{minipage}
		\vspace{1.5em}
	}]
	
	% ============================================================
	% INTRODUÇÃO
	% ============================================================
	\section{Resumo introdutório}
	
	É de interesse por parte deste trabalho apresentar uma análise da integridade estrutural de pisos, bem como o seu devido assentamento, por meio de ensaios não destrutivos baseados em respostas do piso à excitações por impulso. O método experimental consistir-se-á no soltar de uma bolinha de peteca sob uma determinada altura tal que venha a percorrer uma trajetória aproximadamente unidimensional em rota de colisão com uma região do piso pré-selecionada. Por meio desta experimentação, visar-se-á, então, correlacionar a variação da rigidez local com a presença de descontinuidades internas, como vazios e delaminações. 
	
	A partir disto, buscar-se-á registrar, por meio de gravações, sinais acústicos decorrentes dos impactos que serão posteriormente analisados com o uso de um programa feito em MATLAB que realiza transformadas de Fourier de tempo reduzido (TTF) no desejo de estimar a frequência natural amortecida $w_{d}$ [rad/s] do piso. Com isto, fazer-se-á uma modelagem de um sistema massa-mola com estruturação matemática analítica e trivial que descreverá a periodicidade ideal desta bola de peteca. Por meio desta modelagem, comparar-se-á os resultados experimentais com os resultados analíticos.
	
	% ============================================================
	% OBJETIVOS
	% ============================================================
	\section{Objetivos}
	
	Para bem realizar este trabalho, selecionou-se os seguintes objetivos:
	
	\begin{enumerate}[label=\arabic*)]
		\item Determinar uma altura em relação ao piso escolhido.
		\item Escolher uma região do piso para realizar a experimentação.
		\item Ajustar os sensores que serão utilizados.
		\item Iniciar as gravações de vídeo e de áudio.
		\item Soltar a bolinha de peteca sob a altura definida.
		\item Aguardar a bolinha parar de ricochetear.
		\item Encerrar as gravações, bem como guardar os dados coletados.
		\item Realizar os procedimentos anteriores, porém para regiões diferentes do piso.
		\item Utilizar o programa em MATLAB para processar os dados registrados.
		\item Estimar uma vibração livre não amortecida ideal com os valores da experimentação que mais teve oscilações.
	\end{enumerate}
	
	% ============================================================
	% PRÉ-SELECIONAMENTO DE EQUAÇÕES
	% ============================================================
	\section{Pré-selecionamento de equações}
	
Pretende-se, por meio da modelagem analítica, utilizar certas equações, tal como a equação da frequência natural ensinada por \textcite{rao1995mechanical}

\vspace{-0.4cm}
\begin{flalign}
	&w_{n} = \sqrt{\frac{k}{m}}&
\end{flalign}


Onde $w_{n}$ é a frequência natural [rad/s], $k$ é a rigidez [N/m] e $m$ é a massa [kg]. Também utilizar-se-á a equação diferencial do movimento harmônico simples, ensinada também pelo \textcite{rao1995mechanical}

\vspace{-0.4cm}
\begin{flalign}
	&m\frac{d^{2}S(t)}{dt^{2}} +kS(t) = 0&
\end{flalign}

Onde $S(t)$ é a posição em função do tempo [m], tal que a solução para a equação (2) é dada por

\vspace{-0.4cm}
\begin{flalign}
	&S(t) = Asen(w_{n}t+\phi)&
\end{flalign}

Onde $A$ é a amplitude da nossa vibração [m], $t$ é o tempo [s] e $\phi$ é o ângulo  de fase [rad].

	\section{Referências que poderão ser utilizadas}
	
	Para a realização deste trabalho, pretende-se estudar artigos bem consolidados na Literatura. Por hora, conta-se com os trabalhos do \textcite{blaschke2017non}, bem como do \textcite{nikolaev2014estimation} e do \textcite{roebben1997impulse}.

\printbibliography[title={Referências}]
	

	
\end{document}